\section{张量布局、地址流与第一个矩阵向量内核}

\subsection{为什么先讲布局再讲算子优化}

推理引擎最终会把 tokenizer、decoder graph、KV cache、CPU/GPU 后端和性能证据接在一起；但不能从这些大词直接跳到最快 kernel。这里先拿一层矩阵向量乘做显微镜，因为每个 Transformer projection、MLP 子层和 decode 阶段的小 batch 路径，最终都会回到同一组问题：张量在内存里到底怎样排布，热循环每一步访问哪些地址，累加器和输出缓冲怎样被机器指令更新。

本节先放大最小切片中的第一层：

\begin{lstlisting}[numbers=none,caption={本节先追踪这一层}]
hidden = relu(input * W0 + b0)
\end{lstlisting}

为了便于手工检查，先把 batch 固定为 1，把它看成一个矩阵向量乘：

\begin{lstlisting}[numbers=none,caption={矩阵向量乘的 reference 语义}]
for out in 0..N:
  acc = bias[out]
  for k in 0..K:
    acc += input[k] * weight[out, k]
  hidden[out] = max(acc, 0)
\end{lstlisting}

这段 reference 很慢也没关系，它定义了结果。优化版本可以分块、展开、向量化、预打包权重、融合 ReLU，也可以换成 int8 路径；但每一次改造都必须回答同一组问题：\code{input[k]} 是否连续，\code{weight[out,k]} 是否连续，bias 何时加载，acc 用什么宽度保存，输出写到哪里，activation 是分支还是条件数据流。

在 decoder-only Transformer 的 decode 阶段，同样的形状会反复出现。当前 token 的 hidden state 是一条向量，\code{Wq}、\code{Wk}、\code{Wv}、\code{Wo}、\code{Wgate}、\code{Wup}、\code{Wdown} 都是巨大的权重矩阵。单 token decode 时，很多 projection 看起来就是“一个向量乘很多权重行”；只不过 \code{K} 可能是几千，\code{N} 也可能是几千到上万，权重字节和 cache/TLB 压力会远大于教学切片。

\begin{keyidea}
AI 算子的第一层机器事实不是“矩阵乘很重要”，而是：循环次序决定地址流，地址流决定 cache/TLB 和 SIMD 能否工作，累加器宽度决定正确性边界，reference 和 oracle 决定优化是否可信。
\end{keyidea}

\subsection{张量元数据必须落到地址公式}

一个张量视图至少包含 dtype、shape、stride 和 data pointer。以二维权重矩阵 \code{W0[N, K]} 为例，若按行主序连续存储，地址公式是：

\begin{lstlisting}[numbers=none,caption={行主序权重矩阵的地址公式}]
address(W[out, k]) =
  base_W + (out * stride_out + k * stride_k) * sizeof(element)

row-major contiguous:
  stride_out = K
  stride_k   = 1
\end{lstlisting}

若 \code{stride_k == 1}，内层循环扫描 \code{k} 时权重地址连续；若权重来自转置视图，\code{stride_k} 可能变成 \code{N}，内层循环就变成跨步访问。源码里同样写 \code{W[out,k]}，机器看到的地址流却完全不同。第三册后面讲 packing，本质就是把“不适合内核的 stride”复制成“适合热循环的连续块”。

\begin{systemdiagram}{张量视图怎样变成地址流}
\begin{pdfdiagram}[x=1cm,y=1cm]
  \node[book accent node,text width=2.1cm] (meta) at (0,1.0) {TensorView\\dtype\\shape\\stride};
  \node[book teal node,text width=2.1cm] (formula) at (2.75,1.0) {地址公式\\base + index};
  \node[book purple node,text width=2.1cm] (loop) at (5.5,1.0) {循环次序\\out/k};
  \node[book amber node,text width=2.1cm] (stream) at (8.25,1.0) {地址流\\连续/跨步};
  \node[book navy node,text width=2.1cm] (evidence) at (11.0,1.0) {证据\\汇编/测量};

  \draw[book strong arrow] (meta) -- (formula);
  \draw[book strong arrow] (formula) -- (loop);
  \draw[book strong arrow] (loop) -- (stream);
  \draw[book strong arrow] (stream) -- (evidence);

  \node[book label,text width=2.35cm] at (0,-0.55) {不是裸数组，必须带步长和生命周期};
  \node[book label,text width=2.35cm] at (2.75,-0.55) {每个下标都要能展开成字节偏移};
  \node[book label,text width=2.35cm] at (5.5,-0.55) {谁在内层，谁决定高频地址};
  \node[book label,text width=2.35cm] at (8.25,-0.55) {cache line 和 SIMD 只看地址形状};
  \node[book label,text width=2.35cm] at (11.0,-0.55) {用反汇编确认内核真的如此执行};
\end{pdfdiagram}
\diagramnote{本图要证明的命题是：张量布局不是 API 信息，而是热循环地址流的来源。}
\end{systemdiagram}

\subsection{第一个 float32 矩阵向量内核}

先写一个朴素但清楚的 float32 内核：

\begin{lstlisting}[language=C++,caption={float32 reference 风格矩阵向量乘}]
void matvec_relu_f32(float const* input,
                     float const* weight,
                     float const* bias,
                     float* hidden,
                     int n,
                     int k)
{
    for (int out = 0; out < n; ++out) {
        float acc = bias[out];
        float const* row = weight + out * k;
        for (int i = 0; i < k; ++i) {
            acc += input[i] * row[i];
        }
        hidden[out] = acc > 0.0f ? acc : 0.0f;
    }
}
\end{lstlisting}

这段代码的内层地址流很明确：\code{input[i]} 连续读取，\code{row[i]} 连续读取，\code{acc} 通常留在寄存器里，\code{hidden[out]} 每个输出写一次。若 \code{K} 较大，权重流通常主导读取字节；若 \code{N} 较大，同一个 input 向量会被每个输出行复用。后续优化会围绕两件事展开：怎样让 input 和权重的复用更靠近核心，怎样让多个乘加并行进入 SIMD 或多个累加器。

一个抽象汇编账本可以写成：

\begin{lstlisting}[numbers=none,caption={float32 内层循环的典型机器形态}]
load  input[i]        ; stride-1 load
load  weight[row+i]   ; stride-1 load if row-major
mul   input, weight
add   product into acc
branch back to inner loop
\end{lstlisting}

这不是要求某个编译器必须生成标量 \instruction{mul}/\instruction{add}。在合适目标上，优化器可能生成向量 load、FMA、循环展开和水平归约。这里的重点是：无论标量还是向量，内核都必须消耗 input 地址流、weight 地址流和 acc 依赖链；SIMD 只是一次处理多个相邻元素，不能替代布局和地址分析。

\subsection{把内层循环切成基本块}

为了把第一册和第二册的方法接进 AI 算子，不能只停在 C++ 循环。下面这份伪反汇编把同一个 \code{matvec_relu_f32} 切成三个基本块：外层准备一行权重，内层做 dot product，尾部做 ReLU 和写回。

\begin{lstlisting}[numbers=none,caption={matvec 第一层的伪反汇编基本块}]
B0_row:
  acc = load bias[out]
  row = weight + out * K
  i = 0

B1_dot:
  x = load input[i]
  w = load row[i]
  acc = fma_or_mul_add(acc, x, w)
  i = i + 1
  branch i < K back to B1_dot

B2_activate_store:
  y = max(acc, 0)
  store hidden[out] = y
  out = out + 1
  branch out < N back to B0_row
\end{lstlisting}

这份账本故意不写具体寄存器名，因为寄存器分配、展开和向量化会随编译器变化。但它保留了性能分析必须稳定存在的边：\code{B1_dot} 的两条 load 由 \code{i} 产生地址，\code{acc} 是跨内层迭代的归约依赖，\code{i < K} 是内层回边，\code{hidden[out]} 写回发生在归约完成之后。若后续优化把 \code{B1_dot} 展开成四个累加器，变化的是依赖图；若把权重预打包，变化的是 \code{row[i]} 的地址流；若融合下一层，变化的是 \code{B2_activate_store} 的写回边界。

\begin{systemdiagram}{matvec 热区的控制流、地址流和证据闭环}
\begin{pdfdiagram}[x=1cm,y=1cm]
  \node[book accent node,text width=2.25cm] (row) at (0,1.15) {B0 row\\bias + row base};
  \node[book teal node,text width=2.25cm] (dot) at (3.1,1.15) {B1 dot\\load/load/FMA};
  \node[book purple node,text width=2.25cm] (act) at (6.2,1.15) {B2 activate\\max + store};
  \node[book navy node,text width=2.25cm] (oracle) at (9.3,1.15) {oracle\\reference compare};

  \draw[book strong arrow] (row) -- (dot);
  \draw[book strong arrow] (dot) -- (act);
  \draw[book strong arrow] (act) -- (oracle);
  \draw[book back arrow] (dot.south) .. controls +(0,-0.9) and +(0,-0.9) .. node[book label,below=2pt] {i < K} (dot.south);
  \draw[book dashed arrow] (act.north) .. controls +(0,0.95) and +(0,0.95) .. node[book label,above=2pt] {out < N} (row.north);

  \node[book amber node,text width=2.35cm] (input) at (1.55,-1.15) {input 地址\\base + i};
  \node[book amber node,text width=2.35cm] (weight) at (4.65,-1.15) {weight 地址\\row + i};
  \node[book red node,text width=2.35cm] (accdep) at (7.75,-1.15) {acc 依赖\\old -> new};

  \draw[book weak arrow] (dot.south west) -- (input.north);
  \draw[book weak arrow] (dot.south) -- (weight.north);
  \draw[book weak arrow] (dot.south east) -- (accdep.north);

  \node[book label,text width=10.4cm] at (4.65,-2.55) {验证时先确认基本块和回边，再确认 load 地址是否连续，最后用 reference 输出和误差合同证明优化没有改变语义。};
\end{pdfdiagram}
\diagramnote{本图要证明的命题是：AI 内核分析必须同时切出控制流、地址流、归约依赖和正确性证据。}
\end{systemdiagram}

按这张图读后续优化，很多术语会自然归位。循环展开不是“让代码更长”，而是让 \code{B1_dot} 中出现多条相互独立的 acc 链；SIMD 不是“自动快”，而是要求 \code{input} 和 \code{weight} 至少在热路径上形成可装载的相邻 lane；packing 不是“拷贝一份权重”，而是把 \code{weight} 的跨步或间接地址改造成内层连续流；ReLU 融合不是“少一个函数”，而是移动 \code{B2_activate_store} 的边界，减少中间张量的写回和读回。

\subsection{用一张图同时看语义、地址和机器形态}

\begin{systemdiagram}{矩阵向量内核的语义、地址流和硬件账本}
\begin{pdfdiagram}[x=1cm,y=1cm]
  \node[book accent node,text width=2.05cm] (sem) at (0,0) {语义\\dot(input,row)};
  \node[book teal node,text width=2.15cm] (addr0) at (2.8,1.15) {input 流\\base+i};
  \node[book purple node,text width=2.15cm] (addr1) at (2.8,0) {weight 流\\row+i};
  \node[book navy node,text width=2.15cm] (bias) at (2.8,-1.15) {bias[out]\\一次 load};
  \node[book amber node,text width=2.25cm] (acc) at (5.85,0) {accumulator\\寄存器/FMA};
  \node[book red node,text width=2.05cm] (relu) at (8.55,0) {ReLU\\max/compare};
  \node[book teal node,text width=2.1cm] (store) at (11.1,0.72) {hidden[out]\\一次 store};
  \node[book purple node,text width=2.1cm] (proof) at (11.1,-0.72) {证据\\oracle + 汇编};

  \draw[book weak arrow] (sem) -- (addr0);
  \draw[book strong arrow] (sem) -- (addr1);
  \draw[book weak arrow] (sem) -- (bias);
  \draw[book weak arrow] (addr0) -- (acc);
  \draw[book strong arrow] (addr1) -- (acc);
  \draw[book weak arrow] (bias) -- (acc);
  \draw[book strong arrow] (acc) -- (relu);
  \draw[book strong arrow] (relu) -- (store);
  \draw[book weak arrow] (store) -- (proof);

  \node[book label,text width=10.5cm] at (5.45,-2.42) {本图先固定 float32 路径。int8 路径会把 input/weight load 改成窄整数，把 acc 改成 int32，并增加 requantize、round、clamp。};
\end{pdfdiagram}
\diagramnote{本图要证明的命题是：一个 AI 算子要同时有语义账本、地址账本和机器账本，后续优化才能被验证。}
\end{systemdiagram}

按这张图读代码，bias 不是循环外的装饰，它是每个输出通道的一次初始 load；activation 不是数学符号，它要落成 \instruction{max}、条件移动、近似函数或查表路径；输出 store 不是性能主体，但它决定 hidden buffer 的生命周期和后续层的输入布局。若把 activation 与 matvec 融合，可以少一次写出再读回 hidden 的中间流；但融合后 oracle 必须仍能解释输出误差，调试路径也要能定位是哪一层出错。

\subsection{布局错误怎样伪装成算子慢}

常见误判是看到 \code{matvec} 慢，就直接想“需要 SIMD”。但如果权重不是行主序连续，或者每个输出行通过指针数组间接访问，SIMD 也只能在更差的地址流上工作。下面两种形状语义相同，机器成本不同：

\begin{lstlisting}[numbers=none,caption={两种语义相同但地址形状不同的权重访问}]
row-major packed:
  weight[out*K + 0], weight[out*K + 1], weight[out*K + 2] ...

pointer rows:
  rows[out] -> heap block
  rows[out+1] -> another heap block
  row pointers and weight bytes may live on many pages
\end{lstlisting}

第一种形状让内层循环沿着连续权重流前进；第二种形状先追一层指针，再进入某个堆块。若块很小、分配分散或页很多，TLB 和 cache line 利用率都会变差。对于 7B 级推理引擎，模型加载阶段把权重转换为内核友好的 packed layout 往往比在每次推理时容忍任意布局更合理。代价是加载时间、额外内存和格式版本；收益是热路径简单、可测、可向量化。CPU 后端可能选择 cache-friendly panel，GPU 后端可能选择显存连续、coalesced load 或 tensor-core 友好的 tile；两者都必须从同一份张量元数据出发。

\subsection{TensorView 的 C++20 合同}

工程里不能把张量参数写成一堆裸指针和整数。至少需要一个只观察内存、不拥有内存的 view；真正拥有内存的对象可以是 arena、mmap weight、DeviceBuffer 或 aligned buffer。view 的责任是描述地址公式。

\begin{lstlisting}[numbers=none,caption={TensorView 的最小字段}]
TensorView:
  data
  element_type
  rank
  shape[]
  stride[]
  byte_offset
  device
  layout_tag
\end{lstlisting}

它要满足几条不变量：

\begin{itemize}
  \item \code{shape[d] >= 0}，rank 与 shape/stride 长度一致。
  \item view 不拥有内存，不负责释放 \code{data}。
  \item 地址计算必须检查或由 verifier 证明不会越界。
  \item \code{layout_tag} 只是快速路径提示，真正地址仍由 stride 决定。
  \item kernel 若要求 contiguous，必须显式检查，而不是假设 stride 合适。
\end{itemize}

\begin{lstlisting}[numbers=none,caption={contiguous 检查的语义}]
expected = 1
for dim from rank-1 down to 0:
  if stride[dim] != expected:
    return false
  expected *= shape[dim]
return true
\end{lstlisting}

这个检查看起来简单，但能拦住大量错误。比如一个 \code{[T,H]} 视图经过转置后 shape 仍然可能是二维，但 stride 不再是 row-major contiguous；若后端把它送进只支持连续输入的 GEMM kernel，输出就会错或性能极差。正确路线是：verifier 标出非连续，compiler 插入 layout conversion 或选择支持 stride 的 reference/slow path。

\subsection{layout conversion 也要进入 profiler}

很多推理代码把 layout conversion 当成“准备工作”，没有计入耗时。实际它可能非常贵。若每个 decode step 都把某个 activation 从非连续转换成连续，转换成本会进入热路径；若模型加载时一次性把权重 pack 好，成本属于 prepare 阶段，应该和 decode 分开报告。

\begin{lstlisting}[numbers=none,caption={layout conversion event 字段}]
LayoutConversionEvent:
  source_value
  source_layout
  target_layout
  bytes_read
  bytes_written
  stage: prepare | prefill | decode
  reason: backend_requirement | fusion_boundary | debug_oracle
\end{lstlisting}

这能避免一种常见误判：kernel 本身很快，但前后各有一次 layout conversion，端到端反而慢。profiler 若只记录 kernel 时间，就会把瓶颈藏起来。

\subsection{layout 测试矩阵}

张量布局相关测试不能只测连续 happy path。最小测试矩阵应该包含：

\begin{itemize}
  \item row-major contiguous 二维矩阵。
  \item 带 padding 的矩阵，例如每行 stride 大于列数。
  \item 转置 view，shape 正确但 stride 交换。
  \item 子视图，byte offset 非零。
  \item tail shape，\code{K} 或 \code{N} 不能整除 tile。
  \item 量化 packed layout，逻辑坐标与物理位置不同。
\end{itemize}

\begin{lstlisting}[numbers=none,caption={layout oracle 的基本做法}]
for each logical coordinate:
  reference = dense_reference[coordinate]
  actual = read_through_tensor_view(view, coordinate)
  compare actual with reference
\end{lstlisting}

这个测试不是性能测试，而是证明地址公式正确。只要地址公式错，后面 SIMD、GPU、量化都只是更快地读错数据。

\subsection{从 float32 路径过渡到 int8 路径}

量化路径不是把 \code{float} 替换成 \code{int8_t} 那么简单。同一层 matvec 进入 int8 后，机器账本变成：

\begin{lstlisting}[numbers=none,caption={int8 matvec 的机器账本变化}]
load  int8 input values
load  int8 packed weights
convert or widen to int16/int32 lanes
dot product into int32 accumulators
add int32 bias or folded bias
requantize acc to output dtype
apply clamp or activation
store output
\end{lstlisting}

权重字节变小，cache 能容纳更多参数；某些 CPU 有专门的 dot-product 或 madd 指令，能提高吞吐。但累加器必须更宽，重新量化必须处理 scale、舍入和饱和，oracle 也必须从“逐元素几乎相同”变成“误差在合同范围内”。这就是第三册的基本节奏：每个算子先有 float reference，再有地址账本，再有量化语义，最后才进入特定 ISA 和 kernel 优化。

\subsection{把这条账本接回 Transformer}

现在把本节的 \code{matvec} 账本放回 decoder-only Transformer 的一个 decode 步骤。假设当前层输入是一条 hidden 向量 \code{x[hidden]}，这一层至少要做下面几组投影：

\begin{lstlisting}[numbers=none,caption={decode 单 token 中反复出现的投影账本}]
q = matvec(x, Wq)
k = matvec(x, Wk)
v = matvec(x, Wv)
append k, v into KV cache at current position

attn = attention(q, K_cache[0..pos], V_cache[0..pos])
x = x + matvec(attn, Wo)

gate = matvec(norm(x), Wgate)
up   = matvec(norm(x), Wup)
x = x + matvec(silu(gate) * up, Wdown)
\end{lstlisting}

这里的每个 \code{matvec} 都可以用本节的方法拆开：输入向量地址、权重行地址、accumulator 依赖、输出 store、量化 scale 和 oracle。不同的是，真实模型把尺寸和状态放大了。Q/K/V 的输出还要 reshape 成 \code{[heads, head_dim]}；K/V 不是普通临时输出，而要写入第 \code{layer} 层、第 \code{position} 个 token 的 KV cache；attention 会在后续 token 中反复读取这些 cache 行。

\begin{systemdiagram}{从 matvec 切片接回 decoder 层状态}
\begin{pdfdiagram}[x=1cm,y=1cm]
  \node[book accent node,text width=2.15cm] (x) at (0,1.25) {当前 token\\hidden x};
  \node[book teal node,text width=2.15cm] (proj) at (2.75,1.25) {Q/K/V\\matvec};
  \node[book purple node,text width=2.15cm] (cache) at (5.5,1.25) {写 KV cache\\layer,pos};
  \node[book amber node,text width=2.15cm] (attn) at (8.25,1.25) {attention\\读 0..pos};
  \node[book navy node,text width=2.15cm] (mlp) at (11.0,1.25) {MLP\\gate/up/down};

  \draw[book strong arrow] (x) -- (proj);
  \draw[book strong arrow] (proj) -- (cache);
  \draw[book strong arrow] (cache) -- (attn);
  \draw[book strong arrow] (attn) -- (mlp);

  \node[book label,text width=2.4cm] at (0,-0.35) {一条向量，不等于小问题};
  \node[book label,text width=2.4cm] at (2.75,-0.35) {权重行流和 acc 依赖};
  \node[book label,text width=2.4cm] at (5.5,-0.35) {跨 token 状态写入};
  \node[book label,text width=2.4cm] at (8.25,-0.35) {读越来越长的历史};
  \node[book label,text width=2.4cm] at (11.0,-0.35) {最大权重流之一};
\end{pdfdiagram}
\diagramnote{本图要证明的命题是：最小 matvec 账本不是玩具结论，它会直接进入每一层 Transformer decode 的投影、KV cache 和 MLP 路径。}
\end{systemdiagram}

prefill 阶段则把 \code{x} 从一条向量扩成多条 token 行。此时 \code{matvec} 账本会变成 GEMM/GEMV 问题：prompt token 越多，权重复用越明显，packing、tile、SIMD 和 GPU kernel 的价值越容易体现；decode 阶段每次只处理一个或少数 token，权重流和 KV cache 读写又会成为另一种瓶颈。因此性能报告必须把 prefill tokens/s、decode tokens/s 和首 token 延迟分开写。

\subsection{本节验收线}

读完本节，应该能对最小切片的一层 matvec，以及 Transformer 中的 Q/K/V 或 MLP projection，写出下面几件事：

\begin{itemize}
  \item 用 shape 和 stride 推导 \code{input[k]}、\code{weight[out,k]}、\code{bias[out]}、\code{hidden[out]} 和 KV cache 写入位置的地址公式。
  \item 说明行主序连续权重为什么适合内层 \code{k} 循环，转置或指针行为什么会改变地址流。
  \item 把 float32 matvec 写成 load、mul、add、branch、store 的机器账本。
  \item 解释 activation 融合减少了哪条中间内存流，也说明它怎样改变调试和 oracle 边界。
  \item 说明 int8 路径为什么需要 int32 accumulator、requantize 和误差合同。
  \item 区分 prefill 中的多 token GEMM 倾向和 decode 中的单 token GEMV 倾向。
\end{itemize}

这些输出看起来比“会调用矩阵乘 API”低级，但它们是后续 GEMM packing、SIMD、量化校准、算子融合、KV cache 规划、CPU/GPU 后端和推理图内存规划的共同地基。
